\section{AcademicBench}
\label{sec:bench}

\subsection{Desain}
Empat prinsip memandu benchmark ini. (i) Setiap tugas adalah langkah yang dilakukan mahasiswa atau pembimbing saat menulis skripsi. (ii) Skor dihitung secara deterministik di mana pun tugas memungkinkan; juri LLM hanya digunakan untuk kualitas prosa, dan kit menyiapkan lembar penilaian tersamar untuk membandingkan mereka dengan penilai manusia. (iii) Teks skripsi asli hanya digunakan ketika tidak ada alternatif lain, dan tidak pernah dipublikasikan. (iv) Setiap panggilan dicatat: siapa yang menjawab (sebagaimana dilaporkan oleh sistem yang melayaninya), pada reasoning effort berapa, untuk berapa token, dengan biaya berapa. Semua prompt berbahasa Indonesia. Tabel~\ref{tab:tasks} merangkum delapan tugas.

\begin{table*}[t]
  \caption{Delapan tugas AcademicBench. \emph{Pub.} butir publik di repositori; \emph{Priv.} butir privat (hanya hash). Semua butir kecuali cuplikan ekstraksi ditulis untuk benchmark.}
  \label{tab:tasks}
  \centering
  \small
  \begin{tabular}{@{}l p{5.3cm} r r p{6.6cm}@{}}
    \toprule
    Tugas & Apa yang dilakukan model & Pub. & Priv. & Dinilai oleh \\
    \midrule
    \task{structure} & menguraikan subbab pada bab~III (bab metode) dan IV (bab hasil) untuk judul skripsi pada program studi tertentu & 35 & -- & F1 berbobot prevalensi dari keluarga subbab terhadap struktur tipikal program studi \\
    \task{methodology} & menjawab soal pilihan ganda tentang metode penelitian dan statistik (konsep, pemilihan desain atau uji, perhitungan) & 60 & -- & akurasi \\
    \task{extraction} & mengisi catatan metode dengan kosakata tertutup (pendekatan, data, pengumpulan, analisis, perangkat lunak, ukuran sampel, hipotesis) dari cuplikan skripsi yang disamarkan & -- & 30 & kesepakatan bidang dengan label perak lintas-vendor; bukti verbatim \\
    \task{grounded} & menulis tinjauan pustaka 250--400 kata dari enam abstrak yang diberikan, mengutipnya sebagai \texttt{[P1]}--\texttt{[P6]} & 30 & -- & cakupan, kalimat yang dikutip, sitasi tidak valid dan sitasi penulis--tahun, angka yang tidak didukung, panjang; juri LLM \\
    \task{references} & mencantumkan lima rujukan nyata dengan DOI untuk suatu topik penelitian & 20 & -- & pencarian di Crossref (DOI, lalu judul dan penulis pertama); ketepatan DOI; kegagalan diaudit di doi.org, Scopus dan OpenAlex \\
    \task{language} & menulis abstrak dari ringkasan penelitian, atau menulis ulang draf informal ke bahasa Indonesia baku & 30 & -- & panjang, kebocoran bahasa Inggris, kata tidak baku, angka yang dipertahankan dan yang dibuat-buat, istilah kunci, panjang kalimat; juri LLM \\
    \task{coding} & menulis program analisis Python untuk CSV sintetik (regresi, uji asumsi, $t$/Mann--Whitney/Wilcoxon, ANOVA dan uji post-hoc, chi-square, $\alpha$ Cronbach, N-gain, SUS, metrik) & 30 & -- & eksekusi dalam sandbox; statistik yang dicetak dalam toleransi; pass@1 dan lolos setelah satu perbaikan \\
    \task{latex} & memperbaiki dokumen rusak, menyusun tabel dengan booktabs, mengonversi rujukan APA ke BibTeX & 20 & -- & pdflatex dalam sandbox dan teks PDF; perbandingan sel demi sel; isian BibTeX \\
    \bottomrule
  \end{tabular}
\end{table*}

\subsection{Tugas dan butir}
Butir disusun dengan bantuan asisten LLM dan kemudian diverifikasi secara mekanis bila memungkinkan: jawaban yang dihitung dan program referensi dihitung ulang dengan implementasi independen, butir LaTeX dikompilasi, isian BibTeX dan karya sumber dicari di basis data bibliografis, dan sebelum model mana pun dijalankan, kunci jawaban metodologi diaudit oleh peninjau LLM terpisah yang diinstruksikan berperan sebagai penguji. Tinjauan manusia terhadap setiap kunci jawaban merupakan bagian dari validasi yang direncanakan (Bagian~\ref{sec:discussion}).

\paragraph{Struktur} Rujukan untuk sebuah program studi adalah himpunan \emph{keluarga} subbab (mis.\ ``populasi dan sampel'', ``teknik analisis data'', ``gambaran umum lokasi penelitian'', ``pembahasan'') yang muncul di setidaknya 25\,\% skripsi program studi tersebut, beserta prevalensinya, yang diagregasikan dari dataset skripsi lokal satu universitas Indonesia (1{,}251 skripsi; 25 program studi dengan jumlah skripsi yang memadai). Judul-judulnya baru dan fiktif (35 butir; dua judul untuk sepuluh program studi terbesar). Judul subbab model dipetakan ke keluarga yang sama dengan ekspresi reguler yang sama; subbab hasil tanpa judul yang mengikuti pertanyaan penelitian membentuk keluarga tersendiri. Per bab kami menghitung recall, precision, dan F1 berbobot prevalensi; skor butir adalah rata-rata F1 dari kedua bab. Kami juga menandai bab hasil \emph{generik} yang berpusat pada uji asumsi, uji instrumen, atau uji hipotesis pada program studi yang skripsinya tidak memuat uji tersebut.

\paragraph{Metodologi} Enam puluh soal pilihan ganda empat opsi (19 konsep, 25 skenario, 16 komputasi) mencakup statistik deskriptif, validitas dan reliabilitas, skala Likert, pengambilan sampel (Slovin, Krejcie--Morgan), asumsi regresi klasik, regresi, korelasi, uji perbedaan, ANOVA dan percobaan terancang dengan uji post-hoc, chi-square, PLS-SEM, eksperimen kelas dengan N-gain, penelitian tindakan kelas, model R\&D, Miles--Huberman, kriteria trustworthiness, analisis isi dan wacana, penelitian hukum normatif, tinjauan sistematis, pengujian sistem dan SUS, metrik pembelajaran mesin, analisis usaha tani, eksperimen laboratorium dan etika. Jika buku teks berbeda pendapat, rumusan soal menyatakan konvensi yang digunakan. Audit menghasilkan enam revisi: empat butir yang kunci jawabannya tidak unik sebagaimana dirumuskan, dan dua butir yang perumusan atau atribusinya tidak tepat.

\paragraph{Ekstraksi} Tiga puluh cuplikan dari skripsi asli (lima atau empat per fakultas, tujuh fakultas), disamarkan untuk nama, nomor mahasiswa, serta nama dan alamat universitas, memuat judul, pertanyaan penelitian, tujuan, bab metode dan kerangka bab hasil (paling banyak 8{,}700 karakter). Model mengisi catatan JSON dengan kosakata tertutup (14 pendekatan, 12 teknik pengumpulan, 31 teknik analisis, 14 perangkat lunak) dan mengutip bukti. Tidak ada label emas untuk semua butir, sehingga panel yang terdiri dari tiga model kuat dari tiga vendor (gpt-6.1-sol dan claude-sonnet-5 pada reasoning effort rendah, gemini-4-argon) memberikan label perak berdasarkan mayoritas per bidang; sebuah model kandidat dinilai terhadap mayoritas anggota panel \emph{bukan} dari vendornya sendiri (untuk anggota panel: kesepakatan kedua anggota lainnya), dan bidang tanpa kesepakatan tidak dinilai.

\paragraph{Penulisan berlandaskan sumber} Tiga puluh topik yang sering muncul dalam skripsi Indonesia masing-masing dilengkapi dengan enam karya OpenAlex (artikel 2015--2025 dengan abstrak dan setidaknya 20 sitasi), abstrak dipotong pada akhir kalimat sebelum 1{,}200 karakter dan dibekukan dalam sebuah snapshot. Model harus mengutip \texttt{[P$n$]} saja. Skor deterministik merata-rata cakupan sumber, proporsi kalimat dengan sitasi yang valid, tidak adanya sitasi di luar rentang dan sitasi penulis--tahun, proporsi angka yang muncul dalam setidaknya satu sumber, dan kepatuhan terhadap panjang.

\paragraph{Rujukan} Untuk masing-masing dari 20 topik model mencantumkan lima rujukan dengan DOI dan diberi tahu bahwa mencantumkan lebih sedikit lebih baik daripada mengarang. Setiap rujukan dicari di Crossref~\cite{hendricks2020crossref}, lembaga registrasi DOI terbesar untuk DOI jurnal dan prosiding: pertama berdasarkan DOI yang diberikan, kemudian berdasarkan judul dan penulis pertama, dengan aturan pencocokan paper-agent~\cite{arya2026paperagent}: kesamaan judul $\ge 0.90$ (judul dan subjudul Crossref digabung), tahun $\pm 1$ dari salah satu tanggal publikasi karya tersebut, dan nama keluarga penulis pertama; DOI arXiv diperiksa terhadap arXiv. Jawaban Crossref selain ``ditemukan'' atau ``tidak ditemukan'' (batas laju, kesalahan server) membuat butir tidak diberi skor alih-alih dinyatakan salah. Kredit bernilai 1 untuk karya yang ada dengan DOI yang benar, 0.75 tanpa DOI, 0.5 dengan DOI yang salah dan 0 jika karya tidak ditemukan; jumlah kredit dibagi lima, sehingga tidak menjawab bernilai sama dengan mencantumkan rujukan yang tidak dapat diverifikasi. Karena karya nyata bisa saja hilang dari Crossref, sebuah audit mencari ulang setiap rujukan yang tidak ditemukan dengan aturan yang sama---DOI yang diberikan di doi.org (semua lembaga registrasi DOI), judul di Scopus dan di OpenAlex~\cite{priem2022openalex}; Semantic Scholar menolak permintaan kami---dan mengklasifikasikannya sebagai ditemukan, hanya cocok dengan judul karya nyata, atau tidak ditemukan di mana pun. Audit tidak mengubah skor (Bagian~\ref{sec:results}).

\paragraph{Bahasa} Lima belas ringkasan penelitian terstruktur (delapan bidang, rancangan kuantitatif dan kualitatif) harus menjadi satu abstrak 150--250 kata; lima belas draf informal, penuh kata-kata sehari-hari, ejaan tidak baku dan bahasa Inggris yang tidak perlu, harus menjadi bahasa Indonesia baku. Pemeriksaan deterministik menggunakan daftar 438 kata tidak baku beserta bentuk bakunya menurut KBBI/EYD dan daftar kata fungsi bahasa Indonesia dan bahasa Inggris.

\paragraph{Pengodean} Tiga puluh tugas pada data sintetis (seed tetap; berkas CSV adalah data kanonik) menentukan analisis dan pilihannya secara tepat (mis.\ varians sampel, uji dua arah, Student vs.\ Welch, koreksi kontinuitas) dan kunci JSON yang harus dicetak. Nilai rujukan diperiksa silang dengan implementasi independen, dan kesalahan umum (varians populasi, Welch bukannya Student, nilai $p$ yang tidak disesuaikan) berada di luar toleransi. Program yang galat (crash), melewati batas waktu, atau tidak mencetak JSON lengkap mendapat satu giliran perbaikan dengan pesan kesalahan---tidak pernah dengan nilai yang diharapkan.

\paragraph{LaTeX} Delapan dokumen dengan satu hingga tiga kesalahan sisipan (karakter khusus yang tidak di-escape, kurung kurawal yang tidak seimbang, lingkungan yang salah, ekspresi matematika di luar mode matematika, tabel yang cacat) harus dapat dikompilasi dan mempertahankan teksnya; enam tabel harus ditata menggunakan booktabs serta dilengkapi caption dan label; enam rujukan APA (terverifikasi di Crossref atau katalog perpustakaan) harus menjadi entri BibTeX.

\subsection{Juri dan penilai manusia}
Untuk dua tugas prosa, tiga juri dari tiga vendor (model-model panel ekstraksi) menilai setiap jawaban berdasarkan lima kriteria 1--5 (kesetiaan terhadap sumber atau ringkasan, cakupan atau kelengkapan, koherensi, bahasa Indonesia baku, kepatuhan terhadap instruksi) tanpa melihat nama model; sebuah juri tidak pernah menilai model dari vendornya sendiri, sehingga sebuah jawaban menerima dua atau tiga penilaian. Kelunakan setiap juri---rata-rata penyimpangan dari rata-rata juri pada jawaban yang dinilai oleh ketiganya---dikurangi sebelum dirata-rata. Skor butir dari sebuah tugas prosa adalah rata-rata dari skor deterministik dan skor juri. Kami melaporkan kesepakatan juri satu sama lain ($\alpha$ Krippendorff~\cite{hayes2007alpha}) dan dengan pemeriksaan deterministik ($\rho$ Spearman). Kit menghasilkan lembar penilaian tersamar terstratifikasi untuk dua penilai manusia (30 jawaban dinilai dengan rubrik yang sama, 20 catatan ekstraksi label perak untuk diverifikasi) dan menghitung kesepakatan mereka satu sama lain dan dengan para juri; perbandingan ini bukan bagian dari versi saat ini.

\subsection{Sandbox}
Program yang ditulis model dan LaTeX dijalankan dalam namespace Linux user, mount, network, dan PID dengan tmpfs baru di atas direktori home, temporary dan runtime, interpreter terikat dalam mode baca-saja, tanpa jaringan, lingkungan minimal tanpa token, semua capability dihapus sebelum eksekusi, dan batas cgroup pada memori, tugas, dan CPU.
